\appendix

\section{Parity computations} \label{appendix}

In this appendix we record a number of parity computations used throughout the paper. Recall that we use the notation that if $f(x,y)$ is smooth and has a series expansion of the form $f(x,y) \sim \sum_j a_j(y) x^j$, then we denote by $\{f\}_{n}$ the $n^{\mathrm{th}}$ coefficient function $a_n(y)$. Note that with respect to $0$-derivatives (i.e., derivatives with respect to $x \partial_x$ and $x \partial_{y^{\alpha}} = x \partial_{\alpha}$) one may check $\{ x \partial_x f\}_n = n \{f \}_{n},$ and $\{ x \partial_{{\alpha}} f \}_n = \partial_{{\alpha}} \{f\}_{n-1}.$

It is also easy to check the following
\begin{Lemma} \label{lemma:expansion} If $f$ and $g$ are smooth functions which are even as functions of $x$ to order $2j$, then the product $fg$ is even to order $2j$, with $\{ fg\}_0 = \{f\}_0\{g\}_0$ and $\{ fg\}_{2j+1} = \{ f \}_{2j+1} \{g\}_0 + \{f\}_0 \{g\}_{2j+1}$.
\end{Lemma}

Referring to the definition of even metric given on page~\pageref{defn-even} we have
\begin{Lemma} \label{appendix:first-odd-term-olgxx} Let $g$ be a metric which is even to order~$2j$. Suppose that $x$ is a defining function for which
$|{\rm d}x|^2_{\olg} = \olg^{xx} \sim 1 + \sum\limits_{k=1}^{j} A_{2k} x^{2k} + A_{2j+1} x^{2j+1}$; then completing $x$ to a coordinate system $(x,y)$, we have that
\begin{gather*}%\label{eqA.1}
\olg_{xx} = 1 + \sum_{k=1}^{j} a_{2k} x^{2k} + a_{2j+1} x^{2j+1} ,
\end{gather*}
where $a_{2j+1} = -A_{2j+1}$, i.e., $\{ \gbar_{xx} \}_{2j+1} = - \{\gbar^{xx}\}_{2j+1}$.
\end{Lemma}
\begin{proof}We use the formula
\begin{gather*}%\label{eqA.2}
\olg^{xx} \olg_{xx} + \olg^{x\alpha} \olg_{x\alpha} = 1,
\end{gather*}
recalling that since $\olg$ is even to order $2j$, the coefficients $\olg_{\alpha x}$ and $\olg^{\alpha x}$ are odd to order $2j+1$, hence may be ignored since their product do not contribute to an odd term before order~$2j+1$. Writing $\olg_{xx} \sim 1 + \sum\limits_{k=1}^{j} a_{2k} x^{2k} + a_{2j+1} x^{2j+1}$, we see that the coefficient at order~$2j+1$ of the product $\olg^{xx} \olg_{xx}$ equals $a_{2j+1} + A_{2j+1}$. However, this coefficient vanishes, which shows that the term in $\gbar_{xx} $ at order $2j+1$ is $-A_{2j+1}$.
\end{proof}

Now suppose $\olg$ is a metric that is even to order $2j$, i.e., $\olg_{xx}$ and $\olg_{\alpha \beta}$ are even to order $2j$ and $\olg_{x\alpha}$ is odd to order $2j+1$. We record the both the parity and the parity breaking term in the expansions of $0$-derivatives of the various metric components. Note that because of the prefactor $x$ in the derivative, all of these terms vanish at $x=0$; this fact is used frequently in what follows along with Lemma~\ref{lemma:expansion}.
$$
\begin{tabu}{|l|l|l|l|}\hline
\mbox{component} &\mbox{parity to order} &\mbox{coefficient after parity broken} \\
\hline
x \partial_x \olg_{xx} & \mbox{even to order} \ 2j & \{ x \partial_x \olg_{xx} \}_{2j+1} = (2j+1) \{\olg_{xx}\}_{2j+1} \\
x \partial_{\alpha} \olg_{xx} & \mbox{odd to order} \ 2j+1 & \{ x \partial_{\alpha} \olg_{xx} \}_{2j+2} = \partial_{\alpha} \{ \olg_{xx}\}_{2j+1} \\
x \partial_x \olg_{x\mu} & \mbox{odd to order} \ 2j+1 & \{ x \partial_x \olg_{x\mu} \}_{2j+2} = (2j+2) \{\olg_{x\mu}\}_{2j+2} \\
 x \partial_{\nu} \olg_{x\mu} & \mbox{even to order} \ 2j+2 & \{ x \partial_{\nu} \olg_{x\mu} \}_{2j+3} = \partial_{\nu} \{\olg_{x\mu}\}_{2j+2} \\
 x \partial_x \olg_{\alpha \beta} & \mbox{even to order} \ 2j & \{ x \partial_x \olg_{\alpha \beta} \}_{2j+1} = (2j+1) \{ \olg_{\alpha \beta}\}_{2j+1}\\
x \partial_{\nu} \olg_{\alpha \beta}& \mbox{odd to order} \ 2j+1& \{ x \partial_{\nu} \olg_{\alpha\beta} \}_{2j+2} = \partial_{\nu} \{\olg_{\alpha\beta}\}_{2j+1} \\
 \hline
\end{tabu}
$$

\subsection*{Christoffel symbols} We now document the expansion of the Christoffel symbols. Once again each of these vanishes at $x=0$, and it emerges from the computation that any such symbol with an even
number of $x$ components is odd to order $2j+1$. Any symbol with an odd number of $x$ components is even to order $2j$ and \textit{the term breaking parity involves only $xx$ and $\alpha \beta$ components of the metric}. Further, let $\hat\Lambda^{\cdot}_{\cdot \cdot}$ denotes the Christoffel symbols involving only tangential derivatives of tangential components of the metric. The calculations leading to the table below are straightforward but tedious, and any entry with an asterisk is not explicitly needed in the sequel.
$$
\begin{tabu}{|l|l|l|l|}\hline
\mbox{component} &\mbox{parity to order} &\mbox{coefficient after parity broken} \\
\hline
x \olGamma_{xx}^x & \mbox{even to order} \ 2j & \big\{x \olGamma_{xx}^x\big\}_{2j+1} = \frac{1}{2} (2j+1) \{\olg_{xx}\}_{2j+1} \tsep{3pt}\\
 x \olGamma^{x}_{\alpha \beta} & \mbox{even to order} \ 2j & \big\{x \olGamma^{x}_{\alpha \beta} \big\}_{2j+1} = -\frac{1}{2} (2j+1) \{ \olg_{\alpha \beta}\}_{2j+1} \\
x \olGamma_{\alpha x}^x & \mbox{odd to order} \ 2j+1 & \big\{ x \olGamma_{\alpha x}^x \big\}_{2j+2} = \big\{\frac{1}{2}\olg^{xx} x \partial_{\alpha} \olg_{xx}\big\}_{2j+2} + \big\{ \frac{1}{2} \olg^{x\mu} x \partial_x \olg_{\alpha \mu} \big\}_{2j+2} \\
x \olGamma_{\alpha x}^{\gamma} & \mbox{even to order} \ 2j & \big\{x \olGamma_{\alpha x}^{\gamma}\big\}_{2j+1} = \frac{1}{2}(2j+1) h_0^{\gamma \beta} \{ \olg_{\alpha \beta}\}_{2j+1} \\
x \olGamma_{x x}^{\gamma} & \mbox{odd to order} \ 2j+1 & *\\
x \olGamma_{\alpha \beta}^{\gamma} & \mbox{odd to order} \ 2j+1& \big\{x \olGamma_{\alpha \beta}^{\gamma}\big\}_{2j+2} = -\frac{1}{2} \{ \gbar^{\gamma x} \}_1 \{x \partial_x \gbar_{\alpha \beta} \}_{2j+1} + \big\{ x \hat{\Lambda}_{\alpha \beta}^{\gamma}\big\}_{2j+2} \\
 \hline
\end{tabu}
$$

\subsection*{Hessian and Laplacian of $\boldsymbol{x}$}\label{sec:hess-lap-asymp}

Recall that
\begin{gather*}
 [\hess_{\olg} x]_{ij} = \partial_i \partial_j x - \olGamma_{ij}^k \partial_k x. \end{gather*}
Thus, from the previous subsection, since $x$ is a coordinate, both
\begin{gather*} [x \hess_{\olg}(x)]_{xx} = - x \olGamma_{xx}^x \qquad \mbox{and} \qquad [x \hess_{\olg}(x)]_{\alpha \beta} = - x \olGamma_{\alpha \beta}^x
\end{gather*}
are even to order $2j$, and
\begin{gather} \label{eqn:hess-asymp-xx}
\{x \hess_{\olg}(x)_{xx}\}_{2j+1} = -\frac{1}{2} (2j+1) \{\olg_{xx}\}_{2j+1}, \\
\{x \hess_{\olg}(x)_{\alpha \beta}\}_{2j+1} = \frac{1}{2} (2j+1) \{\olg_{\alpha \beta}\}_{2j+1}.\label{eqn:hess-asymp-ab}
\end{gather}
Further,
\begin{gather*}
 [x \hess_{\olg}(x)]_{x\alpha} = -x \olGamma_{x\alpha}^x
\end{gather*}
is odd to order $2j+1$, and
\begin{align}
 \{x \hess_{\olg}(x)_{x\alpha}\}_{2j+2} &= -\left\{\frac{1}{2}\olg^{xx} x \partial_{\alpha} \olg_{xx}\right\}_{2j+2} - \left\{ \frac{1}{2} \olg^{x\mu} x \partial_x \olg_{\alpha \mu} \right\}_{2j+2} \nonumber \\
 &= -\frac{1}{2} \{ x \partial_{\alpha} \olg_{xx} \}_{2j+2} - \frac{1}{2} \{ \olg^{x \mu} \}_1 \{ x \partial_x \olg_{\alpha \mu} \}_{2j+1} \nonumber \\
 &= -\frac{1}{2} \partial_{\alpha} \{\olg_{xx}\}_{2j+1} - \frac{1}{2}(2j+1)\{ \olg^{x \mu} \}_1 \{ \olg_{\alpha \mu}\}_{2j+1},\label{eqn:mixed-ref-1}
\end{align}
where we have used the fact that $x \partial_{\alpha} \olg_{xx}$ vanishes to third order in the second equation above.

We also compute leading components of $x \Delta_{\olg} x \, \gbar_{ij}$. First recall that
\begin{gather*}
x \Delta_{\olg} x = x \gbar^{pq} \hess_{pq} x =-x \gbar^{pq} \olGamma_{pq}^x = -x\gbar^{xx} \olGamma_{xx}^x - x\gbar^{x\alpha} \olGamma_{x\alpha}^x - x\gbar^{\alpha \beta} \olGamma_{\alpha \beta}^x.
\end{gather*}
This expression is even to order $2j+1$, and the middle term is even to order $2j+2$. Thus, using Lemma~\ref{lemma:expansion} and the fact that the product of $x$ with any Christoffel symbol vanishes at $x=0$ we have
\begin{align*}
 \{x \Delta_{\olg} x\}_{2j+1} &= -\{ \gbar^{xx}\}_0 \{ x \olGamma_{xx}^x\}_{2j+1} - \{ \gbar^{\alpha \beta}\}_0 \{ x\olGamma_{\alpha \beta}^x\}_{2j+1} \\
 &= -\frac{1}{2} (2j+1) \{\olg_{xx}\}_{2j+1} +\frac{1}{2} (2j+1) h_0^{\alpha \beta}\{\olg_{\alpha \beta}\}_{2j+1}.
\end{align*}

Assembling the above, we discover
\begin{align} \label{eqn:laplacian-asymp-xx}
\{ x \Delta_{\olg} x \gbar_{xx} \}_{2j+1} &= -\frac{1}{2} (2j+1) \{\olg_{xx}\}_{2j+1} +\frac{1}{2} (2j+1) h_0^{\mu\nu}\{\olg_{\mu\nu}\}_{2j+1}, \\
\{ x \Delta_{\olg} x \gbar_{\alpha \beta} \}_{2j+1} &= \left(-\frac{1}{2} (2j+1) \{\olg_{xx}\}_{2j+1} +\frac{1}{2} (2j+1) h_0^{\mu \nu}\{\olg_{\mu\nu}\}_{2j+1} \right)h^0_{\alpha \beta},\label{eqn:laplacian-asymp-ab}
\end{align}
whereas
\begin{align}
\{ x \Delta_{\olg} x \, \gbar_{x \alpha} \}_{2j+2} &= \{ \olg_{x \alpha}\}_1 \{ x \Delta_{\olg} x\}_{2j+1} \nonumber\\
&= \{ \olg_{x \alpha}\}_1 \left( -\frac{1}{2} (2j+1) \{\olg_{xx}\}_{2j+1} +\frac{1}{2} (2j+1) h_0^{\mu \nu}\{\olg_{\mu \nu}\}_{2j+1}\right). \label{eqn:mixed-ref-2}
\end{align}

\subsection*{The Ricci curvature}

We now perform a similar analysis on the components of the Ricci curvature. We write this out in a bit more detail. First, we have
\begin{gather*}
x^2 \overline{\Rc}_{xx} = x^2 {\overline{Rm}_{\alpha x x}}^{\alpha} = x^2 \partial_{\alpha} \olGamma_{xx}^{\alpha} - x^2 \partial_{x} \olGamma_{\alpha x}^{\alpha} + x\olGamma_{xx}^{s} x \olGamma_{\alpha s}^{\alpha} - x \olGamma_{\alpha x}^{s} x \olGamma_{x s}^{\alpha} \nonumber\\
\hphantom{x^2 \overline{\Rc}_{xx}}{} = x^2 \partial_{\alpha} \olGamma_{xx}^{\alpha} - x^2 \partial_{x} \olGamma_{\alpha x}^{\alpha} + x\olGamma_{xx}^{x} x\olGamma_{\alpha x}^{\alpha} + x\olGamma_{xx}^{\mu} x\olGamma_{\alpha \mu}^{\alpha} - x\olGamma_{\alpha x}^{x} x\olGamma_{x x}^{\alpha} - x\olGamma_{\alpha x}^{\mu} x\olGamma_{x \mu}^{\alpha}\nonumber\\
\hphantom{x^2 \overline{\Rc}_{xx}}{} = x \partial_{\alpha} \big(x \olGamma_{xx}^{\alpha}\big) - x^2 \partial_{x} \olGamma_{\alpha x}^{\alpha} + x\olGamma_{xx}^{x} x\olGamma_{\alpha x}^{\alpha} + x\olGamma_{xx}^{\mu} x\olGamma_{\alpha \mu}^{\alpha} - x\olGamma_{\alpha x}^{x} x\olGamma_{x x}^{\alpha}
 - x\olGamma_{\alpha x}^{\mu} x\olGamma_{x \mu}^{\alpha}\nonumber\\
\hphantom{x^2 \overline{\Rc}_{xx}}{} = x \partial_{\alpha} \big(x \olGamma_{xx}^{\alpha}\big) - x \partial_{x} (x \olGamma_{\alpha x}^{\alpha} ) + x \olGamma_{\alpha x}^{\alpha} + x\olGamma_{xx}^{x} x\olGamma_{\alpha x}^{\alpha} + x\olGamma_{xx}^{\mu} x\olGamma_{\alpha \mu}^{\alpha}\nonumber\\
\hphantom{x^2 \overline{\Rc}_{xx}=}{} - x\olGamma_{\alpha x}^{x} x\olGamma_{x x}^{\alpha} - x\olGamma_{\alpha x}^{\mu} x\olGamma_{x \mu}^{\alpha}.%\label{eqA.4}
\end{gather*}
Comparing the terms in this expression with the tables above, we find $\Rc_{xx}$ is even to order $2j$. Moreover using the parity for the Christoffel symbols already discussed and that each vanishes at $x=0$, the first and the final four terms are even to order $2j+2$ and thus do not contribute to the term at order~$2j+1$. Thus
\begin{gather}
\big\{ x^2 \overline{\Rc}_{xx} \big\}_{2j+1} = \big\{{-} x \partial_{x} \big(x \olGamma_{\alpha x}^{\alpha}\big) + x \olGamma_{\alpha x}^{\alpha} \big\}_{2j+1}
= \frac{1}{2} \big( {-}(2j+1)^2 + 2j+1 \big) h_0^{\alpha \beta} \{ \olg_{\alpha \beta} \}_{2j+1} \nonumber\\
\hphantom{\big\{ x^2 \overline{\Rc}_{xx} \big\}_{2j+1}}{} = -j (2j+1) h_0^{\alpha \beta} \{ \olg_{\alpha \beta} \}_{2j+1}.\label{eqA.5}
\end{gather}

The tangential Ricci components are
\begin{gather*}
x^2 \overline{\Rc}_{\alpha \beta} x^2 {\overline{Rm}_{s \alpha \beta}}^{s}
= x^2 \big( \partial_s \olGamma_{\alpha \beta}^s - \partial_{\alpha} \olGamma_{s \beta}^s + \olGamma_{\alpha \beta}^r \olGamma_{sr}^s - \olGamma_{s \beta}^r \olGamma_{\alpha r}^s \big) \nonumber\\
\hphantom{x^2 \overline{\Rc}_{\alpha \beta} x^2 {\overline{Rm}_{s \alpha \beta}}^{s} }{} = x^2 \big( \partial_x \olGamma_{\alpha \beta}^x + \partial_{\mu} \olGamma_{\alpha \beta}^{\mu} - \partial_{\alpha} \olGamma_{x \beta}^x - \partial_{\alpha} \olGamma_{\mu \beta}^{\mu}
+ \olGamma_{\alpha \beta}^x \olGamma_{xx}^x + \olGamma_{\alpha \beta}^x \olGamma_{\mu x}^{\mu} \nonumber\\
\hphantom{x^2 \overline{\Rc}_{\alpha \beta} x^2 {\overline{Rm}_{s \alpha \beta}}^{s} =}{} + \olGamma_{\alpha \beta}^{\mu} \olGamma_{x\mu}^x + \olGamma_{\alpha \beta}^{\mu} \olGamma_{\lambda \mu}^{\lambda} - \olGamma_{x \beta}^x \olGamma_{\alpha x}^x - \olGamma_{\mu \beta}^x \olGamma_{\alpha x}^{\mu} - \olGamma_{x \beta}^{\mu} \olGamma_{\alpha \mu}^x - \olGamma_{\lambda \beta}^{\mu} \olGamma_{\alpha \mu}^{\lambda} \big).%\label{eqA.6}
\end{gather*}
This entire expression is even to order $2j$, and moreover all terms of the form $\big(x \olGamma\big)\big(x \olGamma\big)$ are even to order $2j+2$, as are the derivative terms involving tangential partial derivatives. We thus find that
\begin{gather}
\big\{ x^2 \overline{\Rc}_{\alpha \beta} \big\}_{2j+1} = \big\{ x^2 \partial_x \olGamma_{\alpha \beta}^x \big\}_{2j+1}
=\big\{ x \partial_x \big( x \olGamma_{\alpha \beta}^x\big) - x \olGamma_{\alpha \beta}^x \big\}_{2j+1}\label{eqA.7}\\
\hphantom{\big\{ x^2 \overline{\Rc}_{\alpha \beta} \big\}_{2j+1}}{}= -\frac{1}{2} (2j+1)^2 \{\olg_{\alpha \beta}\}_{2j+1} + \frac{1}{2} (2j+1) \{\olg_{\alpha \beta}\}_{2j+1}
 = - j (2j+1)\{\olg_{\alpha \beta}\}_{2j+1} .\nonumber
\end{gather}

Now we consider the mixed component $x^2 \overline{\Rc}_{x\alpha}$:
\begin{gather*}
x^2 \overline{\Rc}_{x \alpha} = x^2 {\overline{Rm}_{s x \alpha}}^{s}
= x^2 \big( \partial_s \olGamma_{x \alpha}^s - \partial_{x} \olGamma_{s \alpha}^s + \olGamma_{x \alpha}^r \olGamma_{sr}^s - \olGamma_{s \alpha}^r \olGamma_{x r}^s \big) \\
\hphantom{x^2 \overline{\Rc}_{x \alpha}}{}= x^2 \big( \partial_x \olGamma_{x \alpha}^x + \partial_{\sigma} \olGamma_{x \alpha}^{\sigma} - \partial_{x} \olGamma_{x \alpha}^x - \partial_{x} \olGamma_{\sigma \alpha}^{\sigma}
 + \olGamma_{x \alpha}^x \olGamma_{xx}^x + \olGamma_{x \alpha}^{\nu} \olGamma_{x\nu}^x + \olGamma_{x \alpha}^x \olGamma_{\sigma x}^{\sigma} + \olGamma_{x \alpha}^{\nu} \olGamma_{\sigma \nu}^{\sigma} \\
\hphantom{x^2 \overline{\Rc}_{x \alpha}=}{} - \olGamma_{x \alpha}^x \olGamma_{x x}^x - \olGamma_{x \alpha}^{\nu} \olGamma_{x \nu}^x
- \olGamma_{\sigma \alpha}^x \olGamma_{x x}^{\sigma}- \olGamma_{\sigma \alpha}^{\nu} \olGamma_{x \nu}^{\sigma} \big) \\
\hphantom{x^2 \overline{\Rc}_{x \alpha}}{}= x^2 \big( \partial_{\sigma} \olGamma_{x \alpha}^{\sigma} - \partial_{x} \olGamma_{\sigma \alpha}^{\sigma}+ \olGamma_{x \alpha}^x \olGamma_{\sigma x}^{\sigma} + \olGamma_{x \alpha}^{\nu} \olGamma_{\sigma \nu}^{\sigma} - \olGamma_{\sigma \alpha}^x \olGamma_{x x}^{\sigma}- \olGamma_{\sigma \alpha}^{\nu} \olGamma_{x \nu}^{\sigma} \big).
\end{gather*}
This coefficient is odd to order $2j+1$, thus we will be interested in the coefficient of the first term which breaks parity, at order $2j+2$. Thus purely from parity considerations we find
\begin{gather}
\big\{ x^2\olGamma_{x \alpha}^x \olGamma_{\sigma x}^{\sigma} + x^2\olGamma_{x \alpha}^{\nu} \olGamma_{\sigma \nu}^{\sigma} - x^2 \olGamma_{\sigma \alpha}^x \olGamma_{x x}^{\sigma}- x^2\olGamma_{\sigma \alpha}^{\nu} \olGamma_{x \nu}^{\sigma} \big\}_{2j+2} \nonumber \\
\qquad{} = \big\{x\olGamma_{x \alpha}^x\big\}_1 \big\{ x\olGamma_{\sigma x}^{\sigma}\big\}_{2j+1}+ \big\{x\olGamma_{x \alpha}^{\nu}\big\}_{2j+1} \big\{ x\olGamma_{\sigma \nu}^{\sigma}\big\}_1 \nonumber \\
\qquad\quad {}- \big\{x \olGamma_{\sigma \alpha}^x\big\}_{2j+1} \big\{ x\olGamma_{x x}^{\sigma}\big\}_1 - \big\{x\olGamma_{\sigma \alpha}^{\nu}\big\}_1 \big\{ x\olGamma_{x \nu}^{\sigma}\big\}_{2j+1}.\label{eqn:mixed-ref-3}
\end{gather}
Observe also
\begin{gather}\label{eqn:mixed-ref-4}
\big\{x^2 \partial_{\sigma} \olGamma_{x \alpha}^{\sigma}\big\}_{2j+2} = \big\{ x \partial_{\sigma} \big( x \olGamma_{x \alpha}^{\sigma} \big) \big\}_{2j+2}= \partial_{\sigma}\big\{ x \olGamma_{x \alpha}^{\sigma} \big\}_{2j+1}.
\end{gather}
Finally,
\begin{gather}
\big\{x^2 \partial_{x} \olGamma_{\sigma \alpha}^{\sigma}\big\}_{2j+2} = \big\{x \partial_{x} \big(x\olGamma_{\sigma \alpha}^{\sigma}\big) - x\olGamma_{\sigma \alpha}^{\sigma} \big\}_{2j+2}
 = (2j+1)\big\{x\olGamma_{\sigma \alpha}^{\sigma} \big\}_{2j+2}\nonumber \\
\qquad{} = \frac{-(2j+1)}{2} \{ \gbar^{\sigma x} \}_1 \{x \partial_x \gbar_{\sigma \alpha} \}_{2j+1} + (2j+1)\big\{ x \hat{\Lambda}_{\sigma \alpha}^{\sigma}\big\}_{2j+2} \nonumber\\
\qquad{} = -\frac{(2j+1)^2}{2} \{ \gbar^{\sigma x} \}_1 \{\gbar_{\sigma \alpha} \}_{2j+1} + \frac{1}{2}
\{\gbar^{\sigma \mu}\}_0 \{ x\partial_{\alpha} \gbar_{\mu \sigma}+ x\partial_{\sigma} \gbar_{\alpha \mu} - x\partial_{\mu} \gbar_{\alpha \sigma}\}_{2j+2}.\label{eqn:mixed-ref-5}
\end{gather}

The reader may now verify from equations \eqref{eqn:mixed-ref-3}, \eqref{eqn:mixed-ref-4}, and \eqref{eqn:mixed-ref-5}, that the coefficient $\big\{ x^2 \Rc_{x\alpha} \big\}_{2j+2}$ involves only components $\{ \olg_{xx} \}_{2j+1}$ and $\{ \olg_{\alpha \beta} \}_{2j+1}$ and their tangential derivatives.
